%% kapa: exam=1 round=20 year=2009 period=1 form=A subject=civil_law subject_ko=민법 seq=25 number=25
%% source: Q-Net 감정평가사 기출문제 (cst00309) artlSeq=5024178 file=extracted/1008/5024178/2009년 제20회 감정평가사 1차시험 기출문제/1교시.hwp
%% conversion: kordoc 4.24.51 -> tools/gen1.py (auto); review=proofread-sonnet; flags=none
\begin{question}{25}
표현대리에 관한 설명으로 옳지 \uline{않은} 것은? (다툼이 있으면 판례에 의함)
\begin{choices}
  \item 투자수익보장약정이 강행법규에 위반되어 무효인 경우 표현대리의 법리가 적용될 여지가 없다.
  \item 공법상 행위의 대리권은 표현대리의 기본대리권이 될 수 없다.
  \item 대리권소멸 후의 표현대리가 인정되는 경우에 그 표현대리의 권한을 넘는 대리행위가 있을 때에는 권한을 넘은 표현대리가 성립할 수 있다.
  \item 본인의 직함, 명칭, 상호 등의 사용 허락 또는 묵인도 대리권수여의 표시로 인정될 수 있다.
  \item 유권대리의 주장 속에 표현대리의 주장이 포함된 것으로 볼 수 없다.
\end{choices}
\end{question}
\answer{2}
